% Part: first-order-logic
% Chapter: introduction
% Section: models-theories

\documentclass[../../../include/open-logic-section]{subfiles}

\begin{document}

\olfileid{fol}{int}{mod}

\olsection{মডেল ও তত্ত্ব}

প্রথম-ক্রমের যুক্তিবিদ্যার সংকেতবিন্যাস ও অর্থতত্ত্ব সংজ্ঞায়িত হয়ে গেলে
!!{structure}s ও অর্থগত ধারণাগুলির ধর্ম অনুসন্ধানের কাজ শুরু করা যায়।
!!{derivation} পদ্ধতিও সংজ্ঞায়িত করে সেগুলি নিয়ে অনুসন্ধান করা যায়।
!!{sentence}s-এর কোনো সেটের ক্ষেত্রে প্রশ্ন করা যায়: কোন
!!{structure}s ওই সেটের সব !!{sentence}s-কে সত্য করে? !!{sentence}s-এর
কোনো সেট~$\Gamma$ দেওয়া থাকলে, সেটিকে পরিতৃপ্ত করে এমন
!!a{structure}~$\Struct{M}$-কে \emph{$\Gamma$-র মডেল} বলা হয়। আমরা
$\Gamma$ থেকে শুরু করে তার মডেলগুলি খুঁজতে পারি—সেগুলি কেমন? কত বড় বা
ছোট হওয়া তাদের জন্য আবশ্যক? আবার একটি !!{structure} বা
!!{structure}s-এর একটি সংগ্রহ থেকে শুরু করেও প্রশ্ন করতে পারি: কোন
!!{sentence}s তাদের মধ্যে সত্য? এমন কোনো !!{sentence}s কি আছে, যেগুলি
এই !!{structure}s-কে এই অর্থে \emph{চরিত্রায়িত} করে যে ওই বাক্যগুলি
সেগুলির মধ্যে, এবং শুধু সেগুলির মধ্যেই, সত্য? এই ধরনের প্রশ্ন
\emph{মডেল তত্ত্বের} বিষয়। এগুলিই \emph{স্বতঃসিদ্ধমূলক পদ্ধতির} ভিত্তি:
কোনো তত্ত্বের স্বতঃসিদ্ধ অর্থাৎ !!{sentence}s-এর একটি সেট দিয়ে
!!{structure}s-এর একটি সংগ্রহের বর্ণনা দেওয়া। স্বতঃসিদ্ধগুলি থেকে
প্রথম-ক্রমের যুক্তিবিদ্যায় ঠিক যেসব !!{sentence}s অনুসৃত হয়, সেগুলিই
স্বতঃসিদ্ধগুলির সব মডেলে সত্য—এই পর্যবেক্ষণের ফলে পদ্ধতিটি সম্ভব হয়।

একটি অতি সরল উদাহরণ হিসেবে প্রাক্‌ক্রম বিবেচনা করো। কোনো সেট~$A$-এর উপর
একটি সম্পর্ক~$R$ প্রতিবিম্ব ধর্মবিশিষ্ট ও পরিযায়ী—দুটিই হলে সেটি একটি
প্রাক্‌ক্রম। মাত্র একটি দ্বিপদ সম্পর্ক-সংকেত~$\Obj P$-সহ প্রথম-ক্রমের
ভাষার জন্য !!a{structure} দিতে ঠিক একটি সেট~$A$ এবং তার উপর একটি দ্বিপদ
সম্পর্ক $R \subseteq A \times A$-ই দরকার: আমরা $\Domain{M} = A$ ও
$\Assign{\Obj P}{M} = R$ স্থির করব। $R$ যেহেতু প্রাক্‌ক্রম, তাই সেটি
প্রতিবিম্ব ধর্মবিশিষ্ট ও পরিযায়ী; এই কথাগুলি বলে এমন প্রথম-ক্রমের
যুক্তিবিদ্যার !!{sentence}s-এর একটি সেট~$\Gamma$ পাওয়া যায়:
\begin{align*}
  & \lforall[\Obj v_0][\Atom{\Obj P}{\Obj v_0,\Obj v_0}]\\
  & \lforall[\Obj v_0][\lforall[\Obj v_1][\lforall[\Obj v_2][((\Atom{\Obj P}{\Obj v_0,\Obj v_1} \land \Atom{\Obj P}{\Obj v_1,\Obj v_2}) \lif \Atom{\Obj P}{\Obj v_0,\Obj v_2})]]]
\end{align*}
এই !!{sentence}s শুধু “যেকোনো~$x$-এর জন্য $Rxx$” ($R$ প্রতিবিম্ব
ধর্মবিশিষ্ট) এবং “যখনই $Rxy$ ও $Ryz$, তখন $Rxz$-ও” ($R$ পরিযায়ী)—এই
কথাগুলির সংকেতায়ন। দেখা যায়, !!a{structure}~$\Struct{M}$ এই দুটি
!!{sentence}s-এর সেট~$\Gamma$-র মডেল যদি এবং কেবল যদি $R$ (অর্থাৎ,
$\Assign{\Obj P}{M}$),~$A$-এর (অর্থাৎ, $\Domain{M}$-এর) উপর একটি
প্রাক্‌ক্রম। অন্যভাবে বললে, $\Gamma$-র মডেলগুলি ঠিক প্রাক্‌ক্রমগুলি।
শুধু~$\Obj P$-কে !!{predicate} হিসেবে রাখা প্রথম-ক্রমের ভাষায় প্রকাশযোগ্য
সব প্রাক্‌ক্রমের যেকোনো ধর্ম (যেমন উপরের প্রতিবিম্ব ধর্ম ও পরিযায়িতা)
$\Gamma$-র দুটি !!{sentence}s থেকে অনুসৃত হয় এবং বিপরীতটিও সত্য। তাই
$\Gamma$-র মডেলগুলি সম্বন্ধে যা-ই প্রমাণ করি, তা সব প্রাক্‌ক্রম সম্বন্ধেই
প্রমাণ করা হয়।

যেকোনো বিশেষ তত্ত্ব ও মডেল-শ্রেণির ক্ষেত্রে (যেমন~$\Gamma$ এবং সব
প্রাক্‌ক্রম) সংশ্লিষ্ট প্রথম-ক্রমের ভাষায় কী প্রকাশ করা যায় এবং কী যায়
না—তা নিয়ে আকর্ষণীয় প্রশ্ন থাকবে। আবার কিছু !!{structure}s-এর ধর্ম
সব ভাষা ও মডেল-শ্রেণির ক্ষেত্রেই আকর্ষণীয়; যথা, !!{domain}-এর আকার-সংক্রান্ত
ধর্ম। যেমন, যেকোনো~$n \in \PosInt$-এর জন্য !!{domain}-এ ঠিক
$n$~!!{element}s আছে—এ কথা সব সময় প্রকাশ করা যায়। অসীমসংখ্যক
!!{sentence}s-এর একটি সেট দিয়ে !!{domain} অসীম—এ কথাও প্রকাশ করা যায়।
কিন্তু সংজ্ঞাক্ষেত্র সসীম বা সংজ্ঞাক্ষেত্র !!{nonenumerable}—এ কথা প্রকাশ
করা যায় না। প্রথম-ক্রমের ভাষার সীমাবদ্ধতা সম্বন্ধে এই ফলগুলি সংহতি ও
লোয়েনহাইম--স্কোলেম উপপাদ্যের পরিণাম।

\end{document}
